\begin{figure}[t]
    \centering

    % --- Layout parameters --------------------------------------------
    % Horizontal distance between the centers of two consecutive blocks.
    \pgfmathsetmacro{\blockDistance}{6.6}
    % Width/height of each block (keep in sync with the node styles below).
    \pgfmathsetmacro{\blockSize}{3.6}
    % Fixed text width inside each block, so that longer labels wrap onto
    % extra lines (growing the block downward) instead of growing wider
    % and colliding with the neighboring block.
    \pgfmathsetmacro{\blockTextWidth}{3.1}
    % Text width used for the labels placed above/below the arrows.
    \pgfmathsetmacro{\arrowTextWidth}{2.6}

    % Vertical offset of the legend row below the main diagram.
    \pgfmathsetmacro{\legendYShift}{-3.6}
    % Horizontal offset of the start of the legend row (position of its first icon).
    \pgfmathsetmacro{\legendXShift}{-\blockSize/2 + 0}
    % Horizontal gap between two consecutive legend entries.
    \pgfmathsetmacro{\legendGap}{1.5}
    % Horizontal gap between a legend icon and its own label.
    \pgfmathsetmacro{\legendIconTextGap}{0.3}
    % --------------------------------------------------------------------

\resizebox{\textwidth}{!}{
\begin{tikzpicture}[]
% Nodes
\node[draw, rectangle, align=center, text width=\blockTextWidth cm, minimum width=\blockSize cm, minimum height=\blockSize cm] (box1) at (0, 0) {\hyperref[def:object]{\hl{\textbf{Object}}} \\ the system \\ the observer \\ seeks to understand};
\node[draw, rectangle, right of=box1, align=center, text width=\blockTextWidth cm, minimum width=\blockSize cm, minimum height=\blockSize cm] (box2) at (\blockDistance, 0) {\hyperref[def:explanation]{\hl{\textbf{Explanation}}} \\ \hyperref[def:interpretable]{\hl{interpretable}} \\ information enabling \\ the observer to refine \\ their mental model};
\node[draw, rectangle, right of=box2, align=center, text width=\blockTextWidth cm, minimum width=\blockSize cm, minimum height=\blockSize cm] (box3) at (2*\blockDistance, 0) {\textbf{Mental model} \\ of the observer \\ regarding the object};
% Arrows
\draw[-{Latex[length=3mm]}] (box1) -- node[midway, above, align=center, text width=\arrowTextWidth cm] {Information \\ extraction \\ through the use \\ of an \\ \hyperref[subsec:explainability_method]{\hl{\textbf{explainability}\\\textbf{method}}}} node[midway, below, align=center] {} (box2);
\draw[-{Latex[length=3mm]}] (box2) -- node[midway, above, align=center, text width=\arrowTextWidth cm] {\textbf{Interpretation} \\ by \\ the observer} node[midway, below, align=center] {} (box3);

% Legend: each entry is chained to the previous one via \legendGap /
% \legendIconTextGap, so the whole row reflows consistently if a label
% changes length or if the gaps are retuned.
\node[draw, rectangle, minimum width=0.3cm, minimum height=0.3cm] (legendIcon1) at (\legendXShift, \legendYShift) {};
\node[anchor=west, right=\legendIconTextGap of legendIcon1] (legendText1) {Information structures};

\node[inner sep=0pt, anchor=west, right=\legendGap of legendText1] (legendArrowStart) {};
\coordinate (legendArrowEnd) at ([xshift=0.4cm]legendArrowStart);
\draw[-{Latex[length=1.5mm]}] (legendArrowStart) -- (legendArrowEnd);
\node[anchor=west, right=\legendIconTextGap of legendArrowEnd] (legendText2) {Information transformation};

\node[anchor=west, right=\legendGap of legendText2] (legendIcon3) {\hl{~~}};
\node[anchor=west, right=\legendIconTextGap of legendIcon3] (legendText3) {Concepts discussed in more detail in this section};

% Brace: spans from just past box2 to just past box3, so it automatically
% follows \blockDistance and \blockSize instead of being retuned by hand.
\coordinate (C) at (\blockDistance+\blockSize/2+1, 2.6);
\coordinate (D) at (2*\blockDistance+\blockSize/2+1, 2.6);
\draw [decorate, decoration={brace, amplitude=10pt}, thick]
          (C) -- node[above=12pt, align=center] {Subjective analysis \\ by the \hyperref[def:observer]{\hl{observer}}} (D);
\end{tikzpicture}
    }
\caption{Explanation process in explainable artificial intelligence.}
\label{fig:explanation_process_illustration}
\end{figure}
